\chapter{Expansion in the hypercube graph}
\label{chap:expansion}
\label{sec:expansion}

\begin{definition}[Hypercube graph]\label{def:hypercube-graph}
	\lean{MIPStarRE.LDT.ExpansionHypercubeGraph.rerandomizeCoord, MIPStarRE.LDT.ExpansionHypercubeGraph.rerandomizeCoordWeight}
	\leanok
	The hypercube graph $C$ has vertex set $\F_q^m$ and an edge between points
	that differ in at most one coordinate, including a self-loop at every
	vertex.  A random edge $(u,v)\sim C$ is sampled by drawing $u\sim\F_q^m$,
	$i\sim\{1,\ldots,m\}$, and $x\sim\F_q$ independently and setting $v$ equal
	to $u$ with its $i$-th coordinate replaced by $x$.
\end{definition}

In the remainder of this chapter, $\ket\psi$ is a vector in
$\calH_{\mathrm A}\ot\calH_{\mathrm B}$, not necessarily normalized, and
$A^u$ is Hermitian for every $u\in\F_q^m$.

\begin{definition}[Local and global variance]\label{def:local-and-variance}
	\lean{MIPStarRE.LDT.ExpansionHypercubeGraph.localVariance, MIPStarRE.LDT.ExpansionHypercubeGraph.globalVariance}
	\leanok
	\uses{def:hypercube-graph}
	The local and global variance of $A$ on $\ket\psi$ are
	\begin{align*}
		\mathbf{Var}_{\mathrm{local}}(A,\psi)
		 & =\tfrac12\E_{(u,v)\sim C}\bra\psi(A^u-A^v)^2\ot I\ket\psi,   \\
		\mathbf{Var}_{\mathrm{global}}(A,\psi)
		 & =\tfrac12\E_{u,v\sim\F_q^m}\bra\psi(A^u-A^v)^2\ot I\ket\psi,
	\end{align*}
	with $u,v$ independent in the second line.
\end{definition}

\begin{lemma}[Local-to-global variance]\label{lem:local-to-global}
	\lean{MIPStarRE.LDT.ExpansionHypercubeGraph.localToGlobal, MIPStarRE.LDT.ExpansionHypercubeGraph.matrixLocalToGlobal}
	\leanok
	\uses{def:local-and-variance}
	$\mathbf{Var}_{\mathrm{global}}(A,\psi)\le m\,\mathbf{Var}_{\mathrm{local}}(A,\psi)$.
	Equivalently, for every function $X$ from $\F_q^m$ to a Hilbert space,
	\[
		\E_{u,v\sim\F_q^m}\|X(u)-X(v)\|^2\le m\,\E_{(u,v)\sim C}\|X(u)-X(v)\|^2.
	\]
\end{lemma}
\begin{proof}
	\leanok
	\uses{lem:hypercube-spectral-gap}
	Both variances are quadratic forms of $\ket\psi$ against operators on the
	vertex register tensored with the given operators: the global variance
	uses the projection $I-E$ onto the functions of mean zero, and the local
	variance uses the Laplacian $L$ of $C$.  After normalization the claim is
	the operator inequality $\frac1M(I-E)\le m\,L$, which follows from
	Lemma~\ref{lem:hypercube-spectral-gap}.  The simplified exposition proves
	the same inequality with the product of the coordinate averaging
	projections; the formalization uses the Fourier diagonalization.
\end{proof}

\begin{lemma}[Spectral gap of the hypercube graph]\label{lem:hypercube-spectral-gap}
	\lean{MIPStarRE.LDT.ExpansionHypercubeGraph.hypercubeSpectralGap, MIPStarRE.LDT.ExpansionHypercubeGraph.hypercubeSpectralGap_le_laplacianEigenvalue}
	\leanok
	\uses{def:hypercube-graph}
	Let $M=q^m$, let $K$ be the normalized adjacency operator of $C$ on the
	vertex register, and let $L=\frac1MI-K$ be its Laplacian.  The additive
	characters $\varphi_\alpha$, $\alpha\in\F_q^m$, diagonalize $L$, with
	eigenvalue $\operatorname{wt}(\alpha)/(mM)$, where $\operatorname{wt}(\alpha)$
	is the number of nonzero coordinates of $\alpha$.  In particular every
	nonconstant character has eigenvalue at least $1/(mM)$.
\end{lemma}
\begin{proof}
	\leanok
	Averaging a character over a random edge in direction $i$ multiplies it by
	$0$ or $1$ according as $\alpha_i\ne0$ or $\alpha_i=0$, by orthogonality of
	the characters of $\F_q$.
\end{proof}
